% iterated torus 結び目の refined BPS 不変量と正値性
% コンパイル：cd docs/paper && latexmk -lualatex main.tex
\documentclass[11pt,a4paper]{ltjsarticle}
\usepackage{amsmath,amssymb,amsthm}
\usepackage{booktabs,array,longtable}
\usepackage{pdflscape}
\usepackage[margin=25mm]{geometry}
\usepackage[hidelinks]{hyperref}

\newtheorem{proposition}{命題}[section]
\newtheorem{conjecture}[proposition]{予想}
\theoremstyle{definition}
\newtheorem{definition}[proposition]{定義}
\newtheorem{hypothesis}[proposition]{仮説}
\theoremstyle{remark}
\newtheorem{remark}[proposition]{注意}

\newcommand{\rCS}{\overline{\mathrm{rCS}}}
\newcommand{\ba}{\mathbf{a}}
\newcommand{\bq}{\mathbf{q}}
\newcommand{\bt}{\mathbf{t}}
\newcommand{\hN}{\widehat{N}}
\newcommand{\hf}{\widehat{f}}
\newcommand{\hbN}{\widehat{\mathbf{N}}}

\title{iterated torus 結び目の refined BPS 不変量と正値性\\
\large --- 代数的ケーブルと非代数的ケーブルの比較 ---}
\author{（著者名）}
\date{2026年10月5日}

\begin{document}
\maketitle

\begin{abstract}
Kameyama--Nawata \cite{KN} は，トーラス結び目・トーラス絡み目の refined Chern--Simons（CS）不変量を
resolved conifold 上の refined BPS 不変量 $\hN_{\rho,g,\beta,J_r}$ に書き換え，
それらが非負整数になるという正値性予想を提出した．
本論文では，この枠組みをトーラス結び目のケーブル（iterated torus 結び目）に適用する．
ケーブルの refined 不変量は Cherednik--Danilenko \cite{CD14} の DAHA 構成で定義されている．
我々はその定義を companion の段で線形に分解し，Gorsky--Neguţ \cite{GN} の楕円 Hall 代数の公式で計算した．
対象は三葉結び目 $T(2,3)$ と $T(2,5)$ の 9 個の 2 本ケーブル，彩色は $|\lambda|\le 3$ の全 6 つのヤング図形である．
計算法は二つの方向で検証した．
一つは \cite{KN} の補助ファイルにあるトーラス結び目の値 27 件と表 13 個の完全な再現，
もう一つは \cite{CD14} の DAHA 不変量の明示例 8 件との全係数での一致である．
主な結果は次のとおりである．
refined LMOV 不変量は，小文字の変数 $(a,q,t)$ で求めた $\hf_\rho$ を，読み取る直前に太字の変数へ変換して得た．
判定したのは，$G$ 展開の前の係数 $\hbN_{\rho;i,j,k}$ と，展開の後の refined BPS 不変量 $\hN_{\rho,g,\beta,J_r}$ の両方である．
代数的なケーブル 3 個では，トーラス結び目と同じく，非長方形の $[2,1]$ を含む全彩色で，
$\hbN$ と $\hN$ の両方が同じ符号の下で非負整数になった．
一方，非代数的なケーブル 6 個では，(N2$'$) の正規化の下でいずれも正値性が破れた．
その原因を切り分けると，三つの型に分かれる．
(I) $T(2,3)$ の $(2,3)$, $(2,5)$-ケーブルでは，非彩色の HOMFLY-PT の段階で，偶奇のそろった refinement では正値性が不可能になる．
(II) $T(2,3)$ と $T(2,5)$ の $(2,1)$-ケーブル，$T(2,3)$ の $(2,-1)$-ケーブルでは，HOMFLY-PT は正値性と両立するが，DAHA の不変量 $\rCS_\square$ 自体が負の係数をもつ．
(III) $T(2,3)$ の $(2,7)$-ケーブルの彩色付きの破れは framing に依存し，framing を十分負にずらすと $|\lambda|\le3$ の全彩色で正になる．
トーラス結び目と代数的なケーブルの正値性は，調べた範囲の framing の変化で崩れなかった．
$[1]$ では $\hf_{[1]}$ が rCS と同一なので，ホモロジーとして解釈される不変量は正の整数係数でなければならない．
そこで非代数的なケーブルの $[1]$ について，Khovanov ホモロジーと knot Floer ホモロジーを計算し，Dunfield--Gukov--Rasmussen の手法で superpolynomial の候補を構成した．
DAHA の値はスペクトル系列の必要条件を満たさず，HOMFLY-PT ホモロジーの Poincaré 多項式ではない．
一方，superpolynomial の候補からは，5 結び目すべてで $G$ 展開の前後とも正の refined LMOV 不変量が得られた（$2J_r$ の偶奇は $\mathbf{4_1}$ と同様にそろわない）．
整数性はすべての場合に成り立った．
$(\rho,g,\beta)$ ごとの $2J_r$ の偶奇の一致は，rCS の指数の偶奇の一様性から自動的に従うことを示した．
あわせて，ケーブルの refined 不変量の単項式正規化について，構造的な規則を提案し，検証した．
\end{abstract}

\tableofcontents

%======================================================================
\section{序論}

\subsection{背景}

$S^3$ 上の refined CS 理論 \cite{AS,AS2} は，トーラス結び目に対して
Macdonald 多項式を用いた $S,T$ 行列で定義され，その安定極限は Cherednik の DAHA による不変量 \cite{Ch} と一致する \cite{GN}．
\cite{KN} は large $N$ 双対の refined 版を定式化し，refined CS 不変量 $\rCS_\lambda$ を，
resolved conifold 上の M2--M5 束縛状態の数え上げ $\hN_{\rho,g,\beta,J_r}$ に書き換えた．
そして，トーラス結び目・トーラス絡み目の全ての例で $\hN_{\rho,g,\beta,J_r}\in\mathbb{Z}_{\ge0}$ となることを観察した（正値性予想）．
$\rCS_\lambda$ 自体は，三葉結び目の $[2,1]$ のように非長方形の $\lambda$ で負の係数をもつことが知られている（\cite{KN} とその引用文献を参照）．
にもかかわらず，$\hN$ に書き直すと正になる．
さらに \cite{KN} は，$G$ 展開をする前の $\hf_\rho$ に素朴な変数変換を施した多項式の係数も，非負整数になることを観察した．

非トーラス結び目については，\cite{KN} は彩色 HOMFLY-PT ホモロジーの Poincaré 多項式を入力として，
homologically thin な結び目（$\mathbf{4_1}$ など）を調べた．
しかし，refined 不変量を任意のヤング図形で計算できる非トーラス結び目のクラスは，
現在のところ iterated torus 結び目だけである．
その定義は，Cherednik--Danilenko による DAHA の反復構成 \cite{CD14,CD15} と，
楕円 Hall 代数を用いた Morton--Samuelson の構成 \cite{MS} で与えられている．

\subsection{本論文の内容}

本論文の目的は，\cite{KN} の枠組みを iterated torus 結び目に適用し，
refined BPS 不変量を $|\lambda|\le3$ の全彩色で求め，正値性・整数性を判定することである．
ケーブルの refined 不変量の定義そのものは既知である \cite{CD14}．
本論文で新しく報告するのは，次の (i)--(iv) である．
\begin{enumerate}
\item ケーブル結び目の refined BPS 不変量 $\hN_{\rho,g,\beta,J_r}$ の計算．
  対象は $T(2,3)$ の $(2,q_2)$-ケーブル（$q_2=13,15,1,3,5,7,-1$）と $T(2,5)$ の $(2,21)$, $(2,1)$-ケーブル，彩色は $[1],[2],[1,1],[3],[2,1],[1,1,1]$．
\item 正値性の判定．代数的なケーブルでは全彩色で成り立ち，非代数的なケーブルでは破れる（\S\ref{sec:results}）．
  破れの原因を，HOMFLY-PT による障害，DAHA の不変量自体の負，framing への依存の三つに切り分けた（\S\ref{sec:cause}）．
\item ケーブルの refined 不変量の単項式（framing）正規化の規則の提案と検証（\S\ref{sec:normalization}）．
\item 非代数的なケーブルの $[1]$ での，Khovanov・knot Floer ホモロジーおよび superpolynomial の候補との比較（\S\ref{sec:categorification}）．
\item 計算の過程で見つけた，既存文献の式・規約の不整合の指摘（付録 \ref{app:issues}）．
\end{enumerate}

主張の確からしさを区別するため，計算機で確認したことには「（計算）」と明記し，
文献から直接得られない規則は「仮説」，一般化した主張は「予想」と書く．

\subsection{用語と記法}

本論文では，superpolynomial という語を Gukov らの原典 \cite{DGR,GS} の意味に限って使う．
すなわち，彩色 HOMFLY-PT ホモロジーの Poincaré 多項式として，
微分などの構造的な性質から組み合わせ的に構成される，正の整数係数をもつ多項式である．
本論文で計算する refined CS 不変量 $\rCS_\lambda$（およびそれと一致する DAHA の不変量）はこれとは区別し，
superpolynomial とは呼ばない．

ヤング図形 $\lambda$ の箱 $s=(i,j)$（$i$ 行 $j$ 列，1 始まり）に対し，
腕長 $a(s)=\lambda_i-j$，脚長 $l(s)=\lambda^T_j-i$，co-arm $a'(s)=j-1$，co-leg $l'(s)=i-1$ とする．
$n(\lambda)=\sum_i(i-1)\lambda_i$ と書く．

%======================================================================
\section{準備}

\subsection{変数}

\cite{KN} に従い，$\Omega$ 背景のパラメータ $q=e^{i\epsilon_1}$, $t=e^{i\epsilon_2}$ を用い，$a=t^N$ で $\mathrm{SU}(N)$ 不変量に戻る．
太字の変数（Dunfield--Gukov--Rasmussen \cite{DGR} の $(\ba,\bq,\bt)$）への変換は
\begin{equation}\label{eq:cov}
a=-\ba^2\bt,\qquad q^{1/2}=-\bq\bt,\qquad t^{1/2}=\bq
\end{equation}
である \cite[\S4, 式 (cov)]{KN}．
DAHA の変数 \cite{CD14} とは $(q,t)$ が一致し，$a_{\mathrm{DAHA}}=-a$ である．
例えば三葉結び目の reduced 不変量は $a(1+qt-aq)/\sqrt{qt}$ で，
\eqref{eq:cov} により $\ba^2\bq^{-2}+\ba^2\bq^2\bt^2+\ba^4\bt^3$ \cite{DGR} になり，
DAHA の $1+qt+a_{\mathrm{DAHA}}q$ \cite{CD15} と単項式倍を除いて一致する（計算）．
\cite{DGR} では $T_{2,3}$ は負の交差をもつ三葉結び目であり，本論文もこの向きに従う．

\subsection{Macdonald 多項式と unknot}

$P_\lambda(x;q,t)$ は monic な Macdonald 多項式で，内積は
$\langle p_\lambda,p_\mu\rangle_{q,t}=\delta_{\lambda\mu}z_\lambda\prod_i(1-q^{\lambda_i})/(1-t^{\lambda_i})$ とする．
$g_\lambda=(q/t)^{|\lambda|}/\langle P_\lambda,P_\lambda\rangle_{q,t}$ と置く．
unknot の不変量（unreduced）は
\begin{equation}\label{eq:unknot}
\rCS_\lambda(\bigcirc)=\Big(\frac{t}{a}\Big)^{|\lambda|/2}\prod_{s\in\lambda}\frac{t^{l'(s)}-a\,q^{a'(s)}}{1-q^{a(s)}t^{l(s)+1}}
\end{equation}
で，これは対称な主特殊化 $P_\lambda(t^{\rho})$ である．
\cite{KN} の本文は因子 $(t/a)^{|\lambda|/2}$ のない形を書くが，補助ファイルは \eqref{eq:unknot} を用いており，
\S\ref{sec:bps} の表を再現するにはこちらに合わせる必要がある．
reduced 不変量 $\mathrm{rCS}_\lambda=\rCS_\lambda/\rCS_\lambda(\bigcirc)$ を主に扱う．

\subsection{ケーブルと DAHA の Newton 対}

結び目 $K$ の $(p,q)$-ケーブル $C_{p,q}(K)$ とは，$K$ の管状近傍の境界上の $(p,q)$ 曲線である．
$p$ は $K$ に沿う回数，$q$ は経線方向の回数で，$q$ は $K$ の Seifert framing で測る．
$T(p_1,q_1)$ の $(p_2,q_2)$-ケーブルは，\cite{CD14} の記法では Newton 対
\begin{equation}
\vec r=(p_1,p_2),\qquad \vec s=(q_1,\ q_2-p_1q_1p_2)
\end{equation}
に対応する \cite[\S2.1]{CD14}．
すべての $r_i,s_i>0$ のとき，すなわち $q_2>p_1q_1p_2$ のとき代数的結び目になる．
本論文で扱う結び目を表 \ref{tab:knots} に示す．

\begin{table}[h]
\centering
\caption{扱う結び目．$s_2=q_2-p_1q_1p_2$．}\label{tab:knots}
\begin{tabular}{llccc}
\toprule
ID & 結び目 & $(\vec r,\vec s)$ & 代数的 & \cite{CD14} の明示例 \\
\midrule
K1 & $T(2,3)$ の $(2,13)$-ケーブル & $(2,2),(3,1)$ & ○ & $[1],[2]$ \\
K3 & $T(2,3)$ の $(2,15)$-ケーブル & $(2,2),(3,3)$ & ○ & $[1]$ \\
K6a & $T(2,5)$ の $(2,21)$-ケーブル & $(2,2),(5,1)$ & ○ & --- \\
K2 & $T(2,3)$ の $(2,1)$-ケーブル & $(2,2),(3,-11)$ & × & $[1]$ \\
K4 & $T(2,3)$ の $(2,3)$-ケーブル & $(2,2),(3,-9)$ & × & $[1]$ \\
K5a & $T(2,3)$ の $(2,5)$-ケーブル & $(2,2),(3,-7)$ & × & $[1]$ \\
K5b & $T(2,3)$ の $(2,7)$-ケーブル & $(2,2),(3,-5)$ & × & $[1]$ \\
K5c & $T(2,3)$ の $(2,-1)$-ケーブル & $(2,2),(3,-13)$ & × & $[1]$ \\
K6b & $T(2,5)$ の $(2,1)$-ケーブル & $(2,2),(5,-19)$ & × & --- \\
\bottomrule
\end{tabular}
\end{table}

%======================================================================
\section{計算法}\label{sec:method}

\subsection{楕円 Hall 代数による torus 結び目}

Gorsky--Neguţ \cite{GN} は，安定化した球面 DAHA（楕円 Hall 代数）の元 $P_{kn,km}$ を用いて，
トーラス結び目の DAHA 不変量を
\begin{equation}\label{eq:GN}
\mathcal{P}^\lambda_{n,m}=\varepsilon\big(P_\lambda[p_k\mapsto P_{kn,km}]\cdot1\big)
\end{equation}
と表した．
ここで $P_{k,0}$ は $p_k$ の掛け算，$\varepsilon(f)=f[p_k\mapsto(1-a^k)/(1-t^k)]$ である．
彼らは Hilbert スキームの固定点基底で，$P_{kn,km}$ の行列要素を標準ヤング盤の和として明示的に与えた \cite[\S4]{GN}．
我々はこの式を，各評価点で動的計画法により計算した．

その際，次の二点は \cite{GN} の本文から直接は読み取れず，数値計算で決めた（計算）．
\begin{itemize}
\item \cite[\S4]{GN} の右辺を固定点基底 $I_\mu=\widetilde H_\mu/\nabla_\mu$ の係数と読むと，
  $c_k P_{kn,km}$ に等しくなる．ここで $c_k=(-1)^{k-1}(1-t_w^{-k})/[k]_{q,t_w}$，$t_w=t^{-1}$ である．
  このとき $P_{k,0}$ はちょうど $\varphi(p_k)$ の掛け算になる（$k=1,2,3$ で確認）．
\item この正規化の下で，擬空三角形に対する交換関係 $[P_x,P_y]=\kappa P_{x+y}$ は，
  三つの異なる三角形で同じ定数 $\kappa=(1-q)/t$ をもって成り立つ．
  傾きによらない $\kappa$ は，$SL(2,\mathbb{Z})$ 共変性と整合する．
\end{itemize}
傾きの向きについては，\cite{DGR} の向きの $T(2,2p+1)$ が傾き $(2,-(2p+1))$ に，$T(3,4)$ が $(3,-4)$ に対応する．

\subsection{ケーブルへの拡張}

\cite[(2.12)]{CD14} は，iterated torus 結び目の DAHA-Jones 多項式を次の反復で定義した．
\[
JD_{\vec r,\vec s}(b)=\Big\{\hat\gamma_1\Big(\cdots\big(\hat\gamma_\ell(P_b)/P_b(q^{-\rho_k})\big)\Downarrow\cdots\Big)\Big\}_{ev}
\]
二段の場合，内側の $F_b=\big(\hat\gamma_2(P_b)/P_b(q^{-\rho_k})\big)\Downarrow$ は対称多項式である．
これを $F_b=\sum_\mu d_{b\mu}P_\mu$ と展開すると，$\hat\gamma_1$ は代数の自己同型なので
\begin{equation}\label{eq:cablesum}
JD_{\vec r,\vec s}(b)=\sum_\mu d_{b\mu}\,\big\{\hat\gamma_1(P_\mu)\Downarrow\big\}_{ev}
\end{equation}
が従う（導出）．
すなわち，ケーブルの不変量は，パターン部分の係数 $d_{b\mu}$ と companion の $\mu$ 彩色 unreduced 不変量の和である．

\eqref{eq:GN} の言葉では，次の手順になる．
\begin{enumerate}
\item パターン $F_\lambda=P_\lambda[p_k\mapsto P_{k\cdot(p_2,-s_2)}]\cdot1$ を求める．
\item その冪和表示に $p_k\mapsto P_{k\cdot(p_1,-q_1)}$ を代入して $1$ に作用させる．
\item $\varepsilon$ を取る．
\end{enumerate}

\begin{remark}
研究計画の当初は，非 refined の Rosso--Jones 型ケーブル公式
$W_\lambda(C_{p,q}K)=\theta_\lambda^{-pq}\sum_\mu c^{(p)}_{\lambda\mu}\theta_\mu^{q/p}W_\mu(K)$ を，
$c^{(p)}$ を Macdonald 版プレシズム係数（$q,t$ は変えない）に置き換えて refined 化する方針だった．
しかしこの素朴な持ち上げは，トーラス結び目の段階で既に \cite{KN} の値を再現しない（計算）．
これは，\cite{DMMSS,MMSS} の $\gamma$ 因子や \cite{Dan,Mo18} の報告と整合する．
\eqref{eq:cablesum} の $d_{\lambda\mu}$ は，素朴なプレシズム係数と $T$ 固有値の積ではない．
\end{remark}

\subsection{単項式の正規化}\label{sec:normalization}

DAHA の不変量は $q^\bullet t^\bullet$ 倍を除いてしか定義されない \cite{CD14}．
一方，\cite{KN} の BPS 抽出の結果は，不変量に単項式を掛けると変わってしまう．
そこでケーブルの正規化を次で定める．

\begin{hypothesis}[正規化]\label{hyp:norm}
reduced refined 不変量 $\mathrm{rCS}_\lambda(K)$ は次の二条件で一意に定まる．
\begin{description}
\item[(N1)] $t=q$ での極限が，framing 0 の reduced 彩色 HOMFLY-PT 多項式（\cite{DGR} の規約）に等しい．
\item[(N2$'$)] 単項式は $M_\lambda=M_\square^{|\lambda|}\,T_\lambda^{F}$ の形をもつ．
  ここで $T_\lambda=q^{n(\lambda^T)}t^{-n(\lambda)}$，$F=p_2q_2$ である（トーラス結び目 $T(m,n)$ では $F=mn$）．
  $M_\square$ の $a$ と $q$ の冪と符号は (N1) で決まる．残る $(q/t)$ の冪は，小文字の変数での次数
  $h(a^xq^yt^z)=x+2y$ の最小値が $0$ になるように選ぶ（条件 (N2)）．
\end{description}
\end{hypothesis}

すべての条件は小文字の変数 $(a,q,t)$ の段階で課す．太字の変数は使わない．
(N2) が決めるのは全体の因子 $(q/t)^{c|\lambda|}$（$x$ の尺度変換）だけである．
この因子は refined LMOV 不変量では $2J_r$ を一様に $2c|\rho|$ だけずらすだけで（全体の符号を除く），
整数性と正値性には影響しない．
つまり仮説 \ref{hyp:norm} は，「rCS が全体の因子を除いて \cite{KN} と同じ流儀である」ことの具体的な実現である．

(N2$'$) の形は次の考察から来る．
\cite{KN} の生成関数 $\sum_\lambda\rCS_\lambda g_\lambda P_\lambda(x)$ の plethystic な構造を保つ正規化の自由度は，
$x$ の尺度変換（$c^{|\lambda|}$）と framing の変化（$T_\lambda^k$）だけである．
framing $k$ は $t=q$ で $q^{\kappa_\lambda k}$ として見えるので (N1) で決まり，$c$ は $\square$ で決まる．

仮説 \ref{hyp:norm} は次の (a)--(d) で支持される（いずれも計算）．
\begin{enumerate}
\item[(a)] トーラス結び目 27 件で，GN の値に仮説 \ref{hyp:norm} を課すと，\cite{KN} の補助ファイルの値に全係数で一致する．
\item[(b)] 全ケーブルの全彩色で，ミラー・転置対称性 $\mathrm{rCS}_{\lambda^T}(a,q,t)=\mathrm{rCS}_\lambda(a,t^{-1},q^{-1})$ が単項式まで厳密に成り立つ．
\item[(c)] 全ケーブルで $\hf_\rho$ が Laurent 多項式になる．
\item[(c$'$)] ケーブル用の計算（反復と正規化）で companion を自明な結び目にすると，\cite{KN} の $T(2,3)$, $T(2,5)$ の値が全彩色で全係数一致する．
  したがってケーブルの正規化は，トーラス結び目と同じ流儀である．
\item[(d)] 当初は，各 $\lambda$ に (N1) と (N2)（最低 $\bt$ 次数 $0$）を独立に課していた．
  これはトーラス結び目と K1--K5b では (N2$'$) と同じ結果を与えるが，K5c では $[2]$ で $(q/t)^1$ ずれる．
  そのずれのため，(b) が破れ，$\hf_{[2]}$ に分母 $(1-q^2)(1-qt)$ が残った．
  $[2]$ を $(q/t)^u$ 倍する一パラメータ族の中で，(b) と (c) を同時に満たすのは (N2$'$) の $u$ だけだった．
\end{enumerate}

\subsection{数値計算と厳密化}

すべての量を，有限体 $\mathbb{F}_p$（$p<2^{31}$）上で多数の評価点 $(q,t)$ をまとめて計算した．
$a$ は多項式のまま保つ．
その後，$(q,t)$ について二次元の Newton 補間を行い，二つの素数の中国剰余定理で整数係数を復元した．
復元した多項式は，補間に使っていない点で，三つ目の素数を用いて毎回検証した．
最大の計算（K6a の BPS 不変量，全 6 表）は約 77 分，メモリ 1.1 GB だった．

%======================================================================
\section{refined BPS 不変量の抽出}\label{sec:bps}

\cite{KN} に従い，reformulated invariant $f^q_\mu$, $f^{\bar t}_\mu$ を次で定める．
\begin{align}
\sum_\lambda\rCS_\lambda\,g_\lambda P_\lambda(x;q,t)&=\mathrm{Exp}\Big(\sum_\mu\frac{f^q_\mu}{q^{1/2}-q^{-1/2}}s_\mu(x)\Big),\\
\sum_\lambda\rCS_\lambda\,P_{\lambda^T}(-x;t,q)&=\mathrm{Exp}\Big(\sum_\mu\frac{f^{\bar t}_\mu}{t^{-1/2}-t^{1/2}}s_\mu(x)\Big).
\end{align}
我々はこれを $|\mu|\le3$ で plethystic 対数から一般に導き，\cite[付録 B]{KN} の明示式と照合した．
その際の注意を二つ挙げる．
\begin{itemize}
\item \cite[付録 B]{KN} と補助ファイルの $f^q_\mu$ は，本文の定義から得られる値の $(t/q)^{|\mu|/2}$ 倍である．
  表を作ったのは前者なので，前者に従う．
\item $f^{\bar t}_{[1,1,1]}$ の印刷された式は誤植である（付録 \ref{app:issues}）．
\end{itemize}
$M_{\mu\rho}(t)=\sum_\sigma C_{\mu\sigma\rho}B_\sigma(t)$ を用いて $\hf_\rho=\sum_\mu M^{-1}_{\rho\mu}(t)f^q_\mu$ と定める．
これは
\begin{equation}
\hf_\rho=\sum_{g,\beta,J_r}(-1)^{2J_r}\hN_{\rho,g,\beta,J_r}\,(q^{1/2}-q^{-1/2})^g(t^{-1/2}-t^{1/2})^g\Big(\frac{q}{t}\Big)^{J_r-\beta/2}a^\beta
\end{equation}
の形をもつと予想されている \cite[\S4]{KN}．
$f^q_\mu$, $f^{\bar t}_\mu$, $\hf_\rho$ はすべて小文字の変数 $(a,q,t)$ のまま厳密に求める．
太字の変数への変換は，refined LMOV 不変量を読み取る直前に一度だけ行う．
抽出の手順は補助ファイルに従う．
\begin{enumerate}
\item $|\rho|$ が奇数のとき，$\hf_\rho$ に $a^{\mp1/2}$ を掛ける（$a$ 次数のシフト）．
\item \eqref{eq:cov} で太字の変数に移る．
\item $\ba^{2\beta}$ の係数を $G=(\bq-\bq^{-1})(\bq\bt-\bq^{-1}\bt^{-1})$ の冪で上から順に展開し，
  $\bt^{2J_r}G^g\ba^{2\beta}$ の係数を $\hN_{\rho,g,\beta,J_r}$ と読む．
\end{enumerate}
全体の符号は，$a^{\pm1/2}=(-\ba^2\bt)^{\pm1/2}$ の分枝に由来する規約である．
したがって，正値性は「全体の符号を適当に選べば全項が非負」で判定し，
破れの大きさは「全体の符号を選んだときの負の項の最小数」で測る．
$a$ 次数のシフトの向き（$|\rho|=1$ で $a^{+1/2}$ か $a^{-1/2}$ か）は \cite{KN} の表の中で統一されていない（付録 \ref{app:issues}）．
本論文は $\mathbf{4_1}$ と $T(2,5)$ の表の向き（$|\rho|=1$ で $a^{-1/2}$）に統一する．
この選び方は $\beta$ を一様に付け替えるだけで，正値性の判定には影響しない．

$q^{1/2}$ の符号にも注意が要る．\cite{KN} の本文の \eqref{eq:cov} は $q^{1/2}=-\bq\bt$ だが，
本文に印刷された $T(2,3)$ の例（$\hf_{[1]},\hf_{[2]},\hf_{[1,1]},f^q_{[2]},f^{\bar t}_{[2]}$）は，
トーラス用の補助ファイルの $q^{1/2}=+\bq\bt$ で全体の符号まで再現され，\eqref{eq:cov} では全体の符号だけが反転する（計算）．
項ごとの相対的な符号が $q^{1/2}$ の符号によらないのは，$(A,Q,T)=(a,q,t)^{1/2}$ の指数の偶奇が全項でそろっている場合である．
本論文の rCS と $\hf_\rho$ は，トーラス結び目とケーブルの全てでこの性質をもつ（計算）．
偶奇が混ざるのは $\mathbf{4_1}$ の $\hf_\rho$ だけであり，\cite{KN} の $\mathbf{4_1}$ の表は \eqref{eq:cov} で再現される．
そこで本論文は \eqref{eq:cov} に統一する．

%======================================================================
\section{検証}\label{sec:checks}

計算法と抽出の手順を，次の (1)--(12) で検証した（すべて計算，自動テストとして実装）．
(9)--(12) は，新しく計算したケーブルの rCS が，太字の変数に移る前の小文字の段階で，
全体の因子を除いて \cite{KN} と同じ流儀になっていることの検査である．
\begin{enumerate}
\item Macdonald 多項式の基本性質を確かめた（Schur 極限，直交性とノルム，三角性，主特殊化の閉公式）．
\item トーラス結び目 $T(2,3),T(2,5),T(2,7),T(2,9)$ の $|\lambda|\le3$ 全彩色と $T(3,4)$ の $[1],[2],[1,1]$ の計 27 件で，
  \cite{KN} の補助ファイルの値と全係数で一致した．
  これらの $t=q$ 極限は，Rosso--Jones 公式による彩色 HOMFLY-PT と一致した．
\item \cite[付録 C]{KN} の表 13 個（$T(2,3)$ の 6 個，$\mathbf{4_1}$ の 6 個，$T(2,5)$ の $[1]$）を全項で再現した．
\item 非 refined のケーブル公式（$\theta_\lambda=q^{\kappa_\lambda}a^{|\lambda|/2}$，framing 補正 $\theta_\lambda^{-pq}$）を，
  ケーブルのブレイドから SageMath で計算した HOMFLY-PT と比べ，9 結び目で一致した．
  blackboard framing の補正の向きは，writhe 1 で描いた自明な結び目のケーブルが $T(2,q)$ になることで確かめた．
\item 彩色 HOMFLY-PT について，$H_{\lambda^T}(a,q)=H_\lambda(a,q^{-1})$ と
  $H_\lambda(a,q{=}1)=H_\square(a,1)^{|\lambda|}$ が，9 結び目 $\times$ $|\lambda|\le3$ で成り立った．
\item \cite{CD14} の明示例 8 件と，全係数で一致した（単項式倍を除く）．
  内訳は K1 の $[1],[2]$ と，K2, K3, K4, K5a, K5b, K5c の $[1]$ である．K1 の $[2]$ は 720 項ある．
\item 全評価点で $f^{\bar t}_\mu=\sum_\rho M_{\mu\rho}(q^{-1})\hf_\rho$ が成り立った．
  対象は $T(2,3)$, $T(2,5)$, $\mathbf{4_1}$ と全ケーブルである．
\item An.~Morozov の一次ずれ公式 \cite{AnMo} が，正規化の自由度を除いて成り立った．
  対象は $T(2,3)$, $T(2,5)$, K1, K2 の全彩色である．
\item 偶奇の一様性：全ての reduced rCS（トーラス結び目と全ケーブル）で，$(A,Q,T)$ の指数の偶奇が全項でそろう．
  つまり rCS は，単項式を除いて $a,q,t$ の整数冪の Laurent 多項式である．
\item refined 指数増大性 \cite[\S2]{KN}：$\mathrm{rCS}_{[2]}(q{=}1)=\mathrm{rCS}_{[1]}(q{=}1)^2$，
  $\mathrm{rCS}_{[2,1]}(q{=}1)=\mathrm{rCS}_{[1]}\mathrm{rCS}_{[1,1]}|_{q=1}$，$\mathrm{rCS}_{[3]}(q{=}1)=\mathrm{rCS}_{[1]}(q{=}1)^3$ と，
  その転置（$t=1$）が，トーラス結び目 4 個と全ケーブルで成り立つ．
  これは彩色どうしの単項式の正規化が整合していることの検査になる．
\item 退化したケーブル：ケーブル用の計算で companion を自明な結び目にすると，\cite{KN} の $T(2,3)$, $T(2,5)$ を全彩色で再現する．
\item 変数変換：\cite{KN} の本文に印刷された $T(2,3)$ の例 5 件を，全体の符号まで再現する（\S\ref{sec:bps}）．
\end{enumerate}

%======================================================================
\section{結果}\label{sec:results}

\subsection{正値性・整数性}

refined LMOV 不変量（refined BPS 不変量）は $\hf_\rho$ から読み取る．\cite[\S4]{KN} に従い，次の二つの段階で正の整数係数性を調べた．
いずれも，小文字の変数で求めた $\hf_\rho$ に $a$ 次数のシフトと \eqref{eq:cov} を施した $\ba,\bq,\bt$ の Laurent 多項式から読む．
\begin{itemize}
\item $G$ 展開の前：$\hf_\rho=\sum_{i,j,k}\hbN_{\rho;i,j,k}\,\ba^{2i}\bq^{2j}\bt^k$ の係数 $\hbN_{\rho;i,j,k}$．
  \cite{KN} は，トーラス結び目で全ての $\rho$ について非負整数になることを観察した．
\item $G$ 展開の後：refined BPS 不変量 $\hN_{\rho,g,\beta,J_r}$．
  \cite{KN} の正値性予想の対象であり，トーラス結び目では全ての $\rho$ で非負整数である．
\end{itemize}
結果を表 \ref{tab:before}, \ref{tab:bps} に示す．
すべての場合で $\hf_\rho$ は Laurent 多項式であり，$\hbN$ と $\hN$ はすべて整数だった（計算）．
表 \ref{tab:before} と表 \ref{tab:bps} の両方で \checkmark となる場合（トーラス結び目，$\mathbf{4_1}$，K1, K3, K6a）は，
同じ全体の符号の下で両方が非負になる（計算）．

$f^q_\mu$, $f^{\bar t}_\mu$ そのものの係数は refined LMOV 不変量ではない．
$f^q_\mu=\sum_\rho M_{\mu\rho}(t)\hf_\rho$ は，正の $\hf_\rho$ に符号の混ざった行列 $M(t)$ を掛けたものなので，一般には正にならない．
実際，\cite{KN} の補助ファイルでも $T(2,3)$ の $f^q_{[1,1]}$ は負の係数をもつ．
\cite{KN} が述べた $f^q$, $f^{\bar t}$ の正値性は対称表現 $[r]$ に限った副次的な観察であり，その比較は付録 \ref{app:fqft} にまとめた．

\begin{table}[h]
\centering
\caption{$G$ 展開の前の正値性：$\hbN_{\rho;i,j,k}$（太字の変数での $\hf_\rho$ の係数）．記号は表 \ref{tab:bps} と同じ．}\label{tab:before}
\begin{tabular}{lcccccc}
\toprule
 & $[1]$ & $[2]$ & $[1,1]$ & $[3]$ & $[2,1]$ & $[1,1,1]$ \\
\midrule
$T(2,3)$ & \checkmark & \checkmark & \checkmark & \checkmark & \checkmark & \checkmark \\
$T(2,5)$ & \checkmark & \checkmark & \checkmark & \checkmark & \checkmark & \checkmark \\
$\mathbf{4_1}$ & \checkmark & \checkmark & \checkmark & \checkmark & \checkmark & \checkmark \\
K1 & \checkmark & \checkmark & \checkmark & \checkmark & \checkmark & \checkmark \\
K3 & \checkmark & \checkmark & \checkmark & \checkmark & \checkmark & \checkmark \\
K6a & \checkmark & \checkmark & \checkmark & \checkmark & \checkmark & \checkmark \\
K2 & 10/47 & 152/1136 & 174/1273 & 967/8525 & 1078/9527 & 1156/9990 \\
K4 & 18/42 & 286/1109 & 330/1255 & 2561/8614 & 2850/9646 & 3025/10121 \\
K5a & 8/38 & 409/1226 & 460/1377 & 4441/9758 & 4982/10844 & 5164/11348 \\
K5b & \checkmark & 136/1395 & 149/1537 & 2097/10706 & 2358/11707 & 2351/12202 \\
K5c & 5/57 & 78/1307 & 77/1443 & 400/9501 & 423/10494 & 376/10864 \\
K6b & 30/133 & 576/4400 & 618/4677 & 3626/33214 & 3766/35136 & 3878/36067 \\
\bottomrule
\end{tabular}

\end{table}

\begin{table}[h]
\centering
\caption{$G$ 展開の後の正値性：refined BPS 不変量 $\hN_{\rho,g,\beta,J_r}$．
\checkmark は全体の符号を選べば全項が非負であることを表す．
$m/N$ は，全 $N$ 項のうち，全体の符号を選び直しても $m$ 項が負であることを表す．
$\dagger$ は，$(g,\beta)$ を固定したとき $2J_r$ の偶奇がそろわない組があることを表す．
上の 3 行は \cite{KN} の再現である．}\label{tab:bps}
\begin{tabular}{lcccccc}
\toprule
 & $[1]$ & $[2]$ & $[1,1]$ & $[3]$ & $[2,1]$ & $[1,1,1]$ \\
\midrule
$T(2,3)$ & \checkmark & \checkmark & \checkmark & \checkmark & \checkmark & \checkmark \\
$T(2,5)$ & \checkmark & \checkmark & \checkmark & \checkmark & \checkmark & \checkmark \\
$\mathbf{4_1}$ & \checkmark$^\dagger$ & \checkmark$^\dagger$ & \checkmark$^\dagger$ & \checkmark$^\dagger$ & \checkmark$^\dagger$ & \checkmark$^\dagger$ \\
K1 & \checkmark & \checkmark & \checkmark & \checkmark & \checkmark & \checkmark \\
K3 & \checkmark & \checkmark & \checkmark & \checkmark & \checkmark & \checkmark \\
K6a & \checkmark & \checkmark & \checkmark & \checkmark & \checkmark & \checkmark \\
K2 & 6/48 & 41/1045 & 51/1190 & 253/7852 & 292/8774 & 318/9254 \\
K4 & 20/43 & 195/952 & 224/1090 & 1749/7420 & 1989/8325 & 2128/8795 \\
K5a & 5/47 & 332/1066 & 376/1210 & 3631/8278 & 4066/9235 & 4263/9733 \\
K5b & \checkmark & 66/1305 & 78/1452 & 1060/9612 & 1181/10553 & 1228/11040 \\
K5c & \checkmark & \checkmark & \checkmark & 6/9123 & 7/10053 & \checkmark \\
K6b & 10/178 & 81/4367 & 91/4656 & 512/32175 & 541/34032 & 560/34980 \\
\bottomrule
\end{tabular}

\end{table}

\subsection{例}

非代数的な K2 の $[1]$ の $\hN$ を，各 $(\beta,g)$ について $\sum_{2J_r}\hN\,\bt^{2J_r}$ の形で表 \ref{tab:K2} に示す．
全体の符号は，負の項が少なくなる向きに選んだ．
K1 の $[1]$ は付録 \ref{app:K1} に示す．

\begin{table}[h]
\centering
\caption{K2（$T(2,3)$ の $(2,1)$-ケーブル），$\rho=[1]$ の $\sum_{2J_r}\hN_{[1],g,\beta,J_r}\bt^{2J_r}$．}\label{tab:K2}
\small
\begin{tabular}{c|llll}
\hline
$\beta \backslash g$ & 0 & 1 & 2 & 3 \\
\hline
1 & $-\mathbf{t}^{-1}+\mathbf{t}^{3}+\mathbf{t}^{5}$ & $-2+2\mathbf{t}^{4}$ & $-\mathbf{t}+\mathbf{t}^{3}$ & $0$ \\
2 & $-1+4\mathbf{t}^{4}+5\mathbf{t}^{6}+2\mathbf{t}^{8}$ & $-2\mathbf{t}+2\mathbf{t}^{3}+10\mathbf{t}^{5}+5\mathbf{t}^{7}$ & $-\mathbf{t}^{2}+4\mathbf{t}^{4}+4\mathbf{t}^{6}$ & $\mathbf{t}^{5}$ \\
3 & $5\mathbf{t}^{5}+8\mathbf{t}^{7}+5\mathbf{t}^{9}+\mathbf{t}^{11}$ & $2\mathbf{t}^{4}+13\mathbf{t}^{6}+12\mathbf{t}^{8}+3\mathbf{t}^{10}$ & $3\mathbf{t}^{5}+8\mathbf{t}^{7}+3\mathbf{t}^{9}$ & $\mathbf{t}^{6}+\mathbf{t}^{8}$ \\
4 & $2\mathbf{t}^{6}+5\mathbf{t}^{8}+4\mathbf{t}^{10}+2\mathbf{t}^{12}$ & $5\mathbf{t}^{7}+9\mathbf{t}^{9}+5\mathbf{t}^{11}$ & $4\mathbf{t}^{8}+4\mathbf{t}^{10}$ & $\mathbf{t}^{9}$ \\
5 & $\mathbf{t}^{9}+\mathbf{t}^{11}+\mathbf{t}^{13}$ & $2\mathbf{t}^{10}+2\mathbf{t}^{12}$ & $\mathbf{t}^{11}$ & $0$ \\
\hline
\end{tabular}

\end{table}

\subsection{負の項の位置}

負の項が少ない場合について，少数派の符号の項の位置 $(g,\beta,2J_r)$ と値を表 \ref{tab:minority} に示す．

\begin{table}[h]
\centering
\caption{少数派の符号の項（$(g,\beta,2J_r)\!:\hN$）と，それらの $2J_r-\beta-g$ の値．}\label{tab:minority}
\footnotesize
\begin{tabular}{llp{9.5cm}c}
\toprule
結び目 & $\rho$ & 項 & $2J_r-\beta-g$ \\
\midrule
K2 & $[1]$ & $(0,1,-1)\!:-1$, $(0,2,0)\!:-1$, $(1,1,0)\!:-2$, $(1,2,1)\!:-2$, $(2,1,1)\!:-1$, $(2,2,2)\!:-1$ & $-2$ \\
K5a & $[1]$ & $(0,5,11)\!:-1$, $(0,6,12)\!:-2$, $(0,7,13)\!:-1$, $(1,5,10)\!:-1$, $(1,6,11)\!:-1$ & $4, 6$ \\
K6b & $[1]$ & $(0,3,-1)\!:-1$, $(0,4,0)\!:-1$, $(1,3,0)\!:-4$, $(1,4,1)\!:-4$, $(2,3,1)\!:-6$, $(2,4,2)\!:-6$, $(3,3,2)\!:-4$, $(3,4,3)\!:-4$, $(4,3,3)\!:-1$, $(4,4,4)\!:-1$ & $-4$ \\
K5c & $[3]$ & $(0,2,-1)\!:-2$, $(0,3,0)\!:-6$, $(0,4,1)\!:-6$, $(0,5,2)\!:-2$, $(1,2,0)\!:-4$, $(1,3,1)\!:-5$ & $-3$ \\
K5c & $[2,1]$ & $(0,2,-1)\!:-2$, $(0,3,0)\!:-6$, $(0,4,1)\!:-6$, $(0,5,2)\!:-2$, $(1,2,0)\!:-9$, $(1,3,1)\!:-23$, $(1,4,2)\!:-3$ & $-3$ \\
\bottomrule
\end{tabular}

\end{table}

%======================================================================
\section{$[1]$ での圏化との比較}\label{sec:categorification}

$\hf_{[1]}=(a^{1/2}-a^{-1/2})\,\mathrm{rCS}_{[1]}$（\S\ref{sec:cause}）なので，$[1]$ の refined LMOV 不変量は原理的に rCS そのものである．
したがって $[1]$ では，ホモロジーの Poincaré 多項式として解釈される不変量は，すべて正の整数係数でなければならない．
また \cite{KN} の立場に立つなら，$G$ 展開が閉じて正整数係数になることも，$[1]$ で確かめられなければならない．
本節では，非代数的なケーブル K2, K4, K5a, K5c, K6b の $[1]$ について，本論文の計算（DAHA/楕円 Hall 代数）を，ホモロジーを作る他の構成と比べる．
対照として $T(2,3)$ と代数的なケーブル K1, K3, K6a，および K5b を用いる．

\subsection{圏化の方向性}

非彩色の HOMFLY-PT 多項式の圏化には，次の方向がある．
\begin{enumerate}
\item[(A)] HOMFLY-PT ホモロジー：行列因子化 \cite{KR2} や Soergel 双加群 \cite{Kh07} で定義された三重次数のホモロジー．
\item[(B)] sl$(N)$ ホモロジー（$N=2$ は Khovanov ホモロジー）と knot Floer ホモロジー：二重次数のホモロジー．
  (A) から sl$(N)$ ホモロジーへのスペクトル系列がある \cite{Ras}．
\item[(C)] superpolynomial \cite{DGR}：(A) の Poincaré 多項式を，微分 $d_N$ と対称性などの構造から組み合わせ的に決める構成（予想に基づく）．
\item[(D)] refined CS 不変量・DAHA・楕円 Hall 代数：ホモロジーではない代数的な構成．本論文の計算はこれである．
\item[(E)] 平面曲線特異点の Hilbert スキーム \cite{ORS}：代数的な結び目に限った幾何的な予想．非代数的な結び目には適用できない．
\end{enumerate}
(A) は正のトーラス絡み目について計算されており \cite{HM}，直接計算は 11 交点以下の素な結び目について行われている \cite{NS}．
今回のケーブルのブレイド表示は 13--29 交点であり，その範囲に含まれない．
対象の 5 結び目について，(A) と (C) の計算例は文献に見つからなかった．

\subsection{Khovanov ホモロジーと knot Floer ホモロジー}

ケーブルのブレイド（\S\ref{sec:checks} の (4) で用いたもの）から，
reduced Khovanov ホモロジーを khoca \cite{khoca}（係数 $\mathbb{F}_{101}$）で，knot Floer ホモロジーを SnapPy \cite{snappy} で計算した（計算）．
三葉結び目では，\cite{DGR} の $\mathcal{P}=\ba^2\bq^{-2}+\ba^2\bq^2\bt^2+\ba^4\bt^3$ の $\ba=\bq^2$ 特殊化が Khovanov ホモロジーに一致する．
また $(A,M)=(e_\bq/2,\,e_\bt-e_\ba)$ への移し替えが，knot Floer ホモロジーに一致する．
全ケーブルで，Khovanov ホモロジーは thick だった（表 \ref{tab:homology}）．

\begin{table}[h]
\centering
\caption{reduced Khovanov ホモロジーと knot Floer ホモロジーの総階数．交点数はブレイド表示のもの．L：L-space 結び目．}\label{tab:homology}
\small
\begin{tabular}{lcccccc}
\toprule
 & 交点数 & Kh & $\delta$ の個数 & HFK & L & DAHA の $\rCS_\square$ \\
\midrule
K1 & 19 & 17 & 3 & 13 & ○ & 正 \\
K5b & 13 & 11 & 3 & 7 & ○ & 正 \\
K2 & 17 & 15 & 3 & 7 & × & 負あり \\
K4 & 15 & 13 & 3 & 5 & ○ & 負あり \\
K5a & 13 & 11 & 3 & 5 & ○ & 負あり \\
K5c & 19 & 17 & 3 & 9 & × & 負あり \\
K6b & 29 & 39 & 5 & 15 & × & 負あり \\
\bottomrule
\end{tabular}
\end{table}

Rasmussen のスペクトル系列 \cite{Ras} の微分は，$\ba=\bq^2$ 特殊化の下で $q$ 次数を保ち，$t$ 次数を 1 だけ変える．
したがって，$\mathcal{P}$ が (A) の Poincaré 多項式なら，
\begin{equation}\label{eq:sscond}
\mathcal{P}(\ba=\bq^2)-\mathrm{Kh}=(1+t)\,Q,\qquad Q\ge0
\end{equation}
が成り立つ（必要条件）．
DAHA の $\rCS_\square$ はこの条件を，$T(2,3)$, K1, K3, K6a, K5b では正規化を変えずに満たし，K2, K4, K5a, K5c, K6b では満たさない（計算）．

\subsection{Gukov らの手法による superpolynomial の構成}

\cite{DGR} の手法で，$[1]$ の superpolynomial の候補 $\mathcal{P}(\ba,\bq,\bt)$（正の整数係数）を整数計画で構成した．
総階数を最小にし，次の条件を課した．
\begin{enumerate}
\item[(H)] $\bt=-1$ で HOMFLY-PT 多項式になる．HOMFLY-PT の符号は，Jones 多項式との一致で決める．
\item[($d_{\pm1}$)] 次数 $(-2,2,-1)$, $(-2,-2,-3)$ の微分で 1 項だけ残る．残る項の次数は $(\Sigma,-\Sigma,0)$, $(\Sigma,\Sigma,\Sigma)$ である（$\Sigma$ は S 不変量 \cite{DGR,GGS}）．
\item[($d_2$)] 次数 $(-2,4,-1)$ の微分のホモロジーを $\ba=\bq^2$ に特殊化すると，Khovanov ホモロジーになる．
\item[($d_0$)] 次数 $(-2,0,-3)$ の微分のホモロジーが，knot Floer ホモロジーになる．
\item[(S)] 対称性 $(e_\ba,e_\bq,e_\bt)\leftrightarrow(e_\ba,-e_\bq,e_\bt-e_\bq)$．これは \cite{GHM} が証明した \cite{DGR} の対称性で，\cite{KN} の $q\leftrightarrow t^{-1}$ の対称性の $[1]$ の場合に当たる．
\end{enumerate}
($d_2$), ($d_0$) で各微分が単独でホモロジーを与えることは \cite{DGR} の予想であり，得られるのは候補である．
本論文の rCS は，全結び目で (H) と (S) をずれなしで満たす（計算）．

\begin{table}[h]
\centering
\caption{$[1]$ の superpolynomial の候補と，そこから得た refined LMOV 不変量．前／後は $G$ 展開の前後の負の項の数／項数．$\dagger$：$2J_r$ の偶奇がそろわない．}\label{tab:superpoly}
\small
\begin{tabular}{lcccc}
\toprule
 & 総階数 & DAHA の値との関係 & 前 & 後 \\
\midrule
$T(2,3)$, K1, K5b & 3, 85, 31 & 一致 & 0 & 0 \\
K4 & 27 & 負の項 $-c\,m$ を $+c\,m\,\bt^{-1}$ に置換 & 0/47 & 0/54$^\dagger$ \\
K2 & 41 & 一致しない & 0/59 & 0/63$^\dagger$ \\
K5c & 47 & 大部分が K4 と同じ型 & 0/59 & 0/77$^\dagger$ \\
K5a & 77 & 一致しない（解は 5 個以上） & 0/107 & 0/109$^\dagger$ \\
K6b & 289 & 一致しない & 0/250 & 0/301$^\dagger$ \\
\bottomrule
\end{tabular}
\end{table}

結果を表 \ref{tab:superpoly} に示す（計算）．
\begin{itemize}
\item 対照の $T(2,3)$, K1, K3, K5b では，候補が DAHA の値と全項で一致した．
\item 非代数的な 5 結び目すべてで，正の候補が存在する．
  そこから得た $\hf_{[1]}=(1+\ba^2\bt)\ba^{-2}\bt^{-1}\mathcal{P}$ は，$G$ 展開の前後とも正の整数係数である．
  $2J_r$ の偶奇はそろわず，\cite{KN} の $\mathbf{4_1}$ と同じ型になる．
  これは \S\ref{sec:cause} (2) の「K4, K5a に正の refinement があるなら偶奇はそろわない」と整合する．
\item 一意性：目的関数を乱数の重みに変えて 5 通り解き直すと，K4, K2, K5c では同じ解，K5a では別の解が 5 個得られた（どれも前後とも正）．
  K6b は一意性を調べておらず，探索の格子も狭い．
\item K4 の候補は，DAHA の負の項 $-c\,m$ を $+c\,m\,\bt^{-1}$ に置き換えたものである．
  これは \cite{CD14} が Cab$(5,2)T(3,2)$ の Khovanov 多項式で観察した「$-1\mapsto1/t$」と同じ型である．
\end{itemize}

\subsection{四重次数の色付きホモロジーとの関係}

Gorsky--Gukov--Stošić \cite{GGS} は，長方形の彩色の HOMFLY-PT ホモロジーが四つの次数 $(a,q,t_r,t_c)$ をもつと予想した．
$t_c$ は \cite{AS,DGR} の規約で，refined CS 不変量はこれに対応する．
予想は全ての結び目が対象だが，ケーブルの計算例・言及はない．
$[1]$ では $t_r=t_c$ となり，\cite{DGR} の三重次数に戻る．
したがって，$[1]$ の DAHA の負の係数を $t_r$ と $t_c$ の規約の違いで説明することはできない．
\cite{GGS} が $[1]$ に加える要請のうち，S 不変量の位置は上の ($d_{\pm1}$) に含めた．
新しい微分 $d_{2|1}$, $d_{1|2}$ などは，十分大きな結び目では非彩色でも非自明になりうる．
しかしそのホモロジーの値が未知なので，条件にはしていない．

\subsection{整合性のまとめ}

\begin{table}[h]
\centering
\caption{$[1]$ での DAHA の値（本論文の計算）と他の構成の整合性．}\label{tab:consistency}
\small
\begin{tabular}{lcc}
\toprule
検査 & 代数的なケーブル・K5b & K2, K4, K5a, K5c, K6b \\
\midrule
HOMFLY-PT・Jones・Alexander（Euler 標数） & 整合 & 整合 \\
\cite{DGR} の対称性 & 整合 & 整合 \\
正の整数係数 & 整合 & 不整合 \\
スペクトル系列の必要条件 \eqref{eq:sscond} & 整合 & 不整合 \\
(C) の候補との一致 & 一致 & 不一致 \\
\bottomrule
\end{tabular}
\end{table}

表 \ref{tab:consistency} のとおり，非代数的な 5 結び目の $[1]$ では，DAHA の値は Euler 標数と対称性のレベルでは他の構成と整合する．
しかし，正値性とスペクトル系列のレベルでは整合しない．
したがってこれらの結び目では，DAHA の値は HOMFLY-PT ホモロジーの Poincaré 多項式ではない（計算と \cite{Ras} の定理）．
一方，(C) の候補では $[1]$ の refined LMOV 不変量は $G$ 展開の前後とも正である．
\cite{KN} が非トーラス結び目で superpolynomial を用いたのと同じ枠組みに入れれば，負になるのは (D) を用いた場合に限られる．

%======================================================================
\section{考察}\label{sec:discussion}

\subsection{代数的と非代数的}

代数的なケーブル K1, K3, K6a の refined LMOV 不変量は，トーラス結び目 $T(2,3)$, $T(2,5)$ と同じ振る舞いをした（計算）．
すなわち，$|\lambda|\le3$ の全彩色で，$G$ 展開の前の $\hbN_{\rho;i,j,k}$ と後の $\hN_{\rho,g,\beta,J_r}$ の両方が，同じ符号の下で非負整数になった．
そこで次を予想する．

\begin{conjecture}
代数的な iterated torus 結び目 $K$ と任意のヤング図形 $\rho$ に対し，
仮説 \ref{hyp:norm} で正規化した refined CS 不変量から得られる refined LMOV 不変量について，次が成り立つ．
\begin{enumerate}
\item[(i)] $G$ 展開の前の $\hbN_{\rho;i,j,k}$ と後の $\hN_{\rho,g,\beta,J_r}$ は，同じ全体の符号の下でともに非負整数である．
\item[(ii)] $(\rho,g,\beta)$ を固定すると，非零の $\hN_{\rho,g,\beta,J_r}$ の $2J_r$ の偶奇はそろう．
\end{enumerate}
\end{conjecture}
(ii) は，rCS の指数の偶奇がそろう限り命題 \ref{prop:parity} から自動的に従う．

非代数的なケーブル 6 個は，(N2$'$) の正規化の下で，いずれも少なくとも一つの彩色で正値性が破れた（表 \ref{tab:before}, \ref{tab:bps}）．
K5c（$(2,-1)$-ケーブル）では，展開の後の $\hN$ は $[1],[2],[1,1],[1,1,1]$ で非負なのに，展開の前の $\hbN$ は $[1]$ から負の項をもつ（57 項中 5 項）．
このように非代数的な結び目では，$G$ 展開の前と後で判定が一致しない．
\cite{KN} の正値性予想を調べるには，$\hbN$ と $\hN$ の両方を確かめる必要がある．
破れの原因は \S\ref{sec:cause} で調べる．

\subsection{負の係数の原因}\label{sec:cause}

非代数的なケーブルの負の係数を，次の三つの観点で切り分けた．
計算に用いたスクリプトは付録「データと再現性」に挙げる．
結果を表 \ref{tab:cause} にまとめる．

\paragraph{(1) $[1]$ の負は $\rCS_\square$ そのもの．}
$\hf_{[1]}=(a^{1/2}-a^{-1/2})\,\mathrm{rCS}_{[1]}$ が厳密に成り立つ（計算，K1, K2, K4, K5c）．
この因子は太字の変数では正の因子 $1+\ba^2\bt$ になるので，$[1]$ の負は reduced な rCS 自体の負である．
rCS$_\lambda$ 自体を太字の変数で見ると，トーラス結び目と代数的なケーブルでは $[2,1]$ 以外で非負である．
非代数的なケーブルでは長方形の彩色でも負になる（K5b の $[1]$ だけ非負）（計算）．
これは \cite{CD14} の観察と整合する．\cite{CD14} では Cab$(7,2)T(3,2)$（K5b）は正で，Cab$(5,2)T(3,2)$（K5a）が最初の非正の例である．
$[1]$ では framing は単項式にしか効かないので，この負は正規化によらない．

\paragraph{(2) HOMFLY-PT による障害．}
\begin{proposition}\label{prop:parity}
$\hf_\rho$ の $(A,Q,T)=(a,q,t)^{1/2}$ の指数の偶奇が全項でそろっているとする．
このとき，太字の変数での $\bt$ の次数の偶奇は $(\ba,\bq)$ の次数ごとに決まり，
$G$ 展開の後では $2J_r\equiv\beta+g+c_\rho \pmod 2$（$c_\rho$ は定数）が成り立つ．
\end{proposition}
\begin{proof}
$a^{e_A/2}q^{e_Q/2}t^{e_T/2}$ は \eqref{eq:cov} と $a$ 次数のシフトの後 $\pm\ba^{2m}\bq^{e_Q+e_T}\bt^{m+e_Q}$ になる（$\beta=m$）．
$e_Q,e_T$ の偶奇は一様なので，$\bt$ の次数 $m+e_Q$ の偶奇は $\ba,\bq$ の次数で決まる．
$G=\bq^2\bt-\bt-\bt^{-1}+\bq^{-2}\bt^{-1}$ の各項は $\bt$ の次数の偶奇を一つずらすので，$G^g$ は $g$ だけずらす．
\end{proof}
したがって，$\hbN$（あるいは $\hN$）が全て非負なら，$\bt=-1$（$t=q$）に落とした係数の符号は，この偶奇で決まる $(-1)^k$ に一様にそろう．
これは $t=q$ の値，すなわち彩色 HOMFLY-PT だけで判定できる必要条件である．
K4 と K5a では，非彩色の HOMFLY-PT がこれを破る．
たとえば K4 では，非 refined の $\hN_{\square,0,\beta}$（$\beta=2,\dots,6$）は $-3,2,6,-7,2$ である．
この破れは，Sage でブレイドから計算した HOMFLY-PT でも独立に確かめた（計算，\texttt{tests/test\_negativity\_cause.py}）．
したがって K4, K5a は，$t=q$ で同じ HOMFLY-PT を与え，偶奇がそろうどんな refinement でも正にならない（導出＋計算）．
裏返せば，K4, K5a の彩色 HOMFLY-PT ホモロジーの Poincaré 多項式が存在するなら，それは偶奇がそろわない（\cite{KN} の $\mathbf{4_1}$ と同じ型）．
本論文の rCS（および DAHA の不変量）は偶奇が必ずそろうので，それとは一致できない．
K2, K5c, K6b では，$t=q$ の符号はすべての彩色で正値性と両立する．したがってこれらの負は，DAHA の refinement に固有である．

\paragraph{(3) framing への依存．}
彩色 $\lambda$ の rCS に $T_\lambda^\delta$ を掛けて framing をずらし，$\hf_\rho$ を求め直した（$t=q$ では $\delta=2$ が framing 1 単位）．
$\delta$ が奇数では計算が閉じなかった（$\hf_\rho$ が Laurent 多項式にならないためと考えるが，未確認）．
偶数の $\delta\in[-24,24]$ で調べた結果は次のとおりである（計算）．
\begin{itemize}
\item $T(2,3)$（$n=2$）と K1（$n=2$ は $\delta\in[-16,16]$，$n=3$ は $\delta=\pm8$）：調べた全ての $\delta$ で，$G$ 展開の前後とも正．
\item K5b：$n=2$ は $\delta\le-12$，$n=3$ は $\delta\le-14$ で，$|\lambda|\le3$ の全彩色が前後とも正になる．
  $\delta=0$，すなわち (N1) の framing 0 では負である．
\item K5c：展開の後は，$n=2$ で $\delta\ge0$，$n=3$ で $\delta\ge2$ で正になる．展開の前は $[1]$ の負を引き継ぎ，どの $\delta$ でも負である．
\item K2, K4, K5a, K6b：どの $\delta$ でも負．
\end{itemize}
つまり代数的な結び目の正値性は framing の変化に対して頑健だが，K5b と K5c の彩色付きの破れは framing の選び方に依存する．
K5b を正にする framing に原理的な意味があるかは未解決である．

\begin{table}[h]
\centering
\caption{非代数的なケーブルの負の係数の原因．I：非彩色の HOMFLY-PT による障害．II：DAHA の不変量 $\rCS_\square$ 自体の負（HOMFLY-PT とは両立）．III：彩色付きの破れが framing に依存．}\label{tab:cause}
\begin{tabular}{llll}
\toprule
結び目 & $\rCS_\square$ & 型 & framing を変えたとき \\
\midrule
K4, K5a & 負 & I & 正にならない \\
K2, K6b & 負 & II & 正にならない \\
K5c & 負 & II & 展開の後だけ正になる（$\delta\ge2$） \\
K5b & 非負 & III & $\delta\le-14$ で全彩色が正 \\
\bottomrule
\end{tabular}
\end{table}

\subsection{長方形と非長方形}

refined CS 不変量 $\rCS_\lambda$ 自体が非長方形の彩色で負の係数をもつことは，三葉結び目で既に知られている．
また Cherednik--Philipp \cite{CP} は，DAHA の不変量の正値性について，長方形のヤング図形が例外的であると述べている．
本論文の結果では，代数的なケーブルの $\hbN$ と $\hN$ は，$[2,1]$ でも非負になった（表 \ref{tab:before}, \ref{tab:bps}）．
つまり refined LMOV 不変量の言葉では，長方形の特殊性は消える．
これは \cite{KN} がトーラス結び目で見出した現象が，代数的な iterated torus 結び目でも成り立つことを示している．

\subsection{偶奇と $U(1)_R$}

\cite{KN} は，$(\rho,g,\beta)$ ごとに $2J_r$ の偶奇がそろうことを，
$U(1)_R$ の電荷 $S_R$ が M2--M5 の BPS 状態に自明に作用する証拠とした．
本論文の全ケーブルでも偶奇はそろった（計算）．一方，$\mathbf{4_1}$ ではそろわない \cite{KN}．
ただし命題 \ref{prop:parity} により，この偶奇の一致は rCS の指数の偶奇の一様性から自動的に従う．
rCS が DAHA または楕円 Hall 代数で与えられる限り，偶奇の一致は $S_R$ についての独立な証拠にはならない．
非代数的なケーブルの負の項は，偶奇がそろったまま現れる．
$\hN$ が BPS 状態の空間の次元であるという解釈とは両立しないので，
少なくとも (N2$'$) の正規化では，非代数的なケーブルに \cite{KN} の幾何的な解釈の前提（$\hN=\dim A$）は成り立たないと考えられる．

\subsection{負の項の局在}

表 \ref{tab:minority} の例では，少数派の符号の項は表の端に局在している．
\begin{itemize}
\item K2, K6b, K5c：すべて一本の対角線 $2J_r-\beta-g=\mathrm{const}$ 上にある．
  これは太字の変数で最低 $\bt$ 次数の辺に当たる．
\item K5a：最大の $\beta$ の側にある．
\end{itemize}
破れが「内部」ではなく「境界」に現れることは，負の寄与が少数の余分な状態に由来する可能性を示唆する．
その解釈は今後の課題である．

\subsection{正規化への依存}

一つの彩色の中の相対的な符号は，全項に共通の単項式によらない．
したがって $[1]$ の判定は正規化に依存しない．
一方，$|\rho|\ge2$ の $\hf_\rho$ は彩色どうしの相対的な単項式，特に framing $T_\lambda^\delta$ に依存する（\S\ref{sec:cause} (3)）．
本論文は (N1) によって framing を 0 に固定した．
これはトーラス結び目で \cite{KN} を再現する選び方だが，非代数的なケーブルにも適切かは確かめられていない．
K5b では，framing を負にずらすと正値性が回復する．
仮説 \ref{hyp:norm} を DAHA または楕円 Hall 代数の側から導くこと，とくに非代数的な結び目での framing の選び方は，未解決である．

%======================================================================
\section{結論と今後の課題}

ケーブル結び目の refined BPS 不変量を $|\lambda|\le3$ の全彩色で計算した．
代数的なケーブルでは，refined LMOV 不変量の正値性が，非長方形の彩色まで含めて $G$ 展開の前（$\hbN$）と後（$\hN$）の両方で成り立つ．
これはトーラス結び目と同じである．
非代数的なケーブルでは，(N2$'$) の正規化の下で正値性が破れる．
その原因は結び目によって異なり，HOMFLY-PT の段階の障害（K4, K5a），DAHA の不変量自体の負（K2, K5c, K6b），framing への依存（K5b）に分かれる．
$[1]$ で他の構成と比べると，DAHA の値は Euler 標数と対称性では整合するが，HOMFLY-PT ホモロジーの Poincaré 多項式ではない．
Dunfield--Gukov--Rasmussen の手法による superpolynomial の候補では，非代数的な 5 結び目でも $[1]$ の refined LMOV 不変量が $G$ 展開の前後とも正になる．
今後の課題を挙げる．
\begin{itemize}
\item 3 本ケーブル（$T(2,3)$ の $(3,19)$-ケーブルなど）と，3 段以上の反復．
\item 単項式の正規化の理論的な導出．
\item 非代数的な結び目で $\hN$ が負になることの物理的解釈．
\item 非代数的な結び目での framing の選び方（K5b で正値性を回復する framing の意味）．
\item HOMFLY-PT ホモロジーの直接計算（\cite{NS} のプログラム）による superpolynomial の候補の確認．
\item 候補の一意性（K5a, K6b）と，\cite{GGS} の新しい微分 $d_{n|m}$ の影響．
\item DAHA の値を superpolynomial から得る規則（\cite{CD14} の「$-1\mapsto1/t$」の一般化）．
\item monotone link・Coxeter link \cite{GL} への拡張．
\end{itemize}

%======================================================================
\appendix
\section{既存文献で見つけた不整合}\label{app:issues}

\begin{enumerate}
\item \cite[付録 B]{KN} の $f^{\bar t}_{[1,1,1]}$ の印刷された式は，補助ファイルの式（\texttt{ft111}）および我々の導出と一致しない．後者二つは一致する．
\item \cite{KN} の本文の $g_\lambda$ による定義と，付録 B の $f^q_\mu$ は，尺度 $(q/t)^{|\mu|/2}$ だけ異なる．
\item \cite[付録 C]{KN} の表の $a$ 次数シフトは統一されていない．
  $T(2,3)$ の表は $|\rho|=1$ で $a^{+1/2}$，$|\rho|=3$ で $a^{-1/2}$ を用い，$T(2,5)$ と $\mathbf{4_1}$ の表は逆を用いる（補助ファイルで確認）．
\item \cite{KN} の本文は $q^{1/2}=-\bq\bt$ だが，トーラス用の補助ファイルの関数 \texttt{knotparam} は $q^{1/2}=+\bq\bt$ である．
  $T(2,3)$ の表は，どちらでも絶対値が一致する．
\item \cite[\S4]{GN} の行列要素の公式を楕円 Hall 代数の元として使うための正規化（\S\ref{sec:method}）．
\end{enumerate}

\section{補足：$f^q_\mu$, $f^{\bar t}_\mu$ の係数}\label{app:fqft}

\cite[\S4]{KN} は副次的な観察として，トーラス結び目の対称表現 $\mu=[r]$ では
$f^q_{[r]}$, $f^{\bar t}_{[r]}$ を \eqref{eq:cov} で太字の変数に移すと，全体の $\pm\bq^\bullet$ を除いて非負整数係数 $\mathbf{N}^q_{[r];i,j,k}$, $\mathbf{N}^{\bar t}_{[r];i,j,k}$ になると述べた．
これらは refined LMOV 不変量ではなく，$[r]$ 以外では正になるとは主張されていない．
比較のため，表 \ref{tab:Nq}, \ref{tab:Nt} に全彩色の結果を示す．
$T(2,3)$ の $[1,1],[2,1],[1,1,1]$ の負の項の数は，\cite{KN} の補助ファイルで同じ式を評価した結果と一致する（計算）．
代数的なケーブル K1, K3, K6a でも，対称表現 $[1],[2],[3]$ では非負になった（計算）．
$\mathbf{N}^q$ と $\mathbf{N}^{\bar t}$ の負の項の数は全ての場合で一致した．
これは自由エネルギーの関係 $F^{\bar t}(a,q,t)=F^q(a,t^{-1},q^{-1})$ \cite{KN} を反映している．

\begin{table}[h]
\centering
\caption{$\mathbf{N}^q_{\mu;i,j,k}$（$f^q_\mu$ の係数）．refined LMOV 不変量ではない．記号は表 \ref{tab:bps} と同じ．}\label{tab:Nq}
\begin{tabular}{lcccccc}
\toprule
 & $[1]$ & $[2]$ & $[1,1]$ & $[3]$ & $[2,1]$ & $[1,1,1]$ \\
\midrule
$T(2,3)$ & \checkmark & \checkmark & 17/51 & \checkmark & 81/217 & 105/340 \\
$T(2,5)$ & \checkmark & \checkmark & 65/162 & \checkmark & 308/735 & 354/1216 \\
$\mathbf{4_1}$ & \checkmark & 19/38 & 35/70 & 22/180 & 128/320 & 144/452 \\
K1 & \checkmark & \checkmark & 1767/4043 & \checkmark & 9062/20864 & 8701/34023 \\
K3 & \checkmark & \checkmark & 2250/5006 & \checkmark & 11415/25879 & 10791/42190 \\
K6a & \checkmark & \checkmark & 5774/11936 & \checkmark & 29302/60871 & 24736/98502 \\
K2 & 10/47 & 122/686 & 466/1331 & 675/3684 & 2077/6035 & 3022/10586 \\
K4 & 18/42 & 210/651 & 515/1253 & 1411/3609 & 2320/6202 & 4093/10425 \\
K5a & 8/38 & 232/662 & 654/1308 & 1497/3883 & 3019/6930 & 5378/11258 \\
K5b & \checkmark & 35/790 & 732/1501 & 520/4495 & 3259/7315 & 4483/11989 \\
K5c & 5/57 & 59/771 & 484/1527 & 293/4094 & 2320/6544 & 2728/11661 \\
K6b & 30/133 & 440/2514 & 1696/4764 & 2301/13372 & 7339/20518 & 9484/37151 \\
\bottomrule
\end{tabular}

\end{table}

\begin{table}[h]
\centering
\caption{$\mathbf{N}^{\bar t}_{\mu;i,j,k}$（$f^{\bar t}_\mu$ の係数）．refined LMOV 不変量ではない．}\label{tab:Nt}
\begin{tabular}{lcccccc}
\toprule
 & $[1]$ & $[2]$ & $[1,1]$ & $[3]$ & $[2,1]$ & $[1,1,1]$ \\
\midrule
$T(2,3)$ & \checkmark & \checkmark & 17/51 & \checkmark & 81/217 & 105/340 \\
$T(2,5)$ & \checkmark & \checkmark & 65/162 & \checkmark & 308/735 & 354/1216 \\
$\mathbf{4_1}$ & \checkmark & 19/38 & 35/70 & 22/180 & 128/320 & 144/452 \\
K1 & \checkmark & \checkmark & 1767/4043 & \checkmark & 9062/20864 & 8701/34023 \\
K3 & \checkmark & \checkmark & 2250/5006 & \checkmark & 11415/25879 & 10791/42190 \\
K6a & \checkmark & \checkmark & 5774/11936 & \checkmark & 29302/60871 & 24736/98502 \\
K2 & 10/47 & 122/686 & 466/1331 & 675/3684 & 2077/6035 & 3022/10586 \\
K4 & 18/42 & 210/651 & 515/1253 & 1411/3609 & 2320/6202 & 4093/10425 \\
K5a & 8/38 & 232/662 & 654/1308 & 1497/3883 & 3019/6930 & 5378/11258 \\
K5b & \checkmark & 35/790 & 732/1501 & 520/4495 & 3259/7315 & 4483/11989 \\
K5c & 5/57 & 59/771 & 484/1527 & 293/4094 & 2320/6544 & 2728/11661 \\
K6b & 30/133 & 440/2514 & 1696/4764 & 2301/13372 & 7339/20518 & 9484/37151 \\
\bottomrule
\end{tabular}

\end{table}

\begin{landscape}
\section{K1 の $\rho=[1]$}\label{app:K1}
\begin{center}
\tiny
\begin{tabular}{c|>{\raggedright\arraybackslash}p{4.0cm}>{\raggedright\arraybackslash}p{4.0cm}>{\raggedright\arraybackslash}p{4.0cm}>{\raggedright\arraybackslash}p{4.0cm}>{\raggedright\arraybackslash}p{4.0cm}}
\hline
$g \backslash \beta$ & 7 & 8 & 9 & 10 & 11 \\
\hline
0 & $\mathbf{t}^{-1}+\mathbf{t}+2\mathbf{t}^{3}+3\mathbf{t}^{5}+4\mathbf{t}^{7}+4\mathbf{t}^{9}+4\mathbf{t}^{11}+3\mathbf{t}^{13}+\mathbf{t}^{15}$ & $1+2\mathbf{t}^{2}+4\mathbf{t}^{4}+7\mathbf{t}^{6}+10\mathbf{t}^{8}+12\mathbf{t}^{10}+12\mathbf{t}^{12}+10\mathbf{t}^{14}+4\mathbf{t}^{16}$ & $\mathbf{t}^{3}+2\mathbf{t}^{5}+5\mathbf{t}^{7}+8\mathbf{t}^{9}+12\mathbf{t}^{11}+13\mathbf{t}^{13}+12\mathbf{t}^{15}+6\mathbf{t}^{17}$ & $\mathbf{t}^{8}+2\mathbf{t}^{10}+4\mathbf{t}^{12}+6\mathbf{t}^{14}+6\mathbf{t}^{16}+4\mathbf{t}^{18}$ & $\mathbf{t}^{15}+\mathbf{t}^{17}+\mathbf{t}^{19}$ \\
1 & $8+14\mathbf{t}^{2}+23\mathbf{t}^{4}+32\mathbf{t}^{6}+37\mathbf{t}^{8}+34\mathbf{t}^{10}+23\mathbf{t}^{12}+8\mathbf{t}^{14}$ & $8\mathbf{t}+21\mathbf{t}^{3}+41\mathbf{t}^{5}+66\mathbf{t}^{7}+85\mathbf{t}^{9}+87\mathbf{t}^{11}+65\mathbf{t}^{13}+26\mathbf{t}^{15}$ & $7\mathbf{t}^{4}+18\mathbf{t}^{6}+39\mathbf{t}^{8}+60\mathbf{t}^{10}+73\mathbf{t}^{12}+63\mathbf{t}^{14}+30\mathbf{t}^{16}$ & $5\mathbf{t}^{9}+12\mathbf{t}^{11}+20\mathbf{t}^{13}+23\mathbf{t}^{15}+14\mathbf{t}^{17}$ & $2\mathbf{t}^{16}+2\mathbf{t}^{18}$ \\
2 & $28\mathbf{t}+63\mathbf{t}^{3}+100\mathbf{t}^{5}+124\mathbf{t}^{7}+118\mathbf{t}^{9}+79\mathbf{t}^{11}+28\mathbf{t}^{13}$ & $28\mathbf{t}^{2}+84\mathbf{t}^{4}+160\mathbf{t}^{6}+230\mathbf{t}^{8}+247\mathbf{t}^{10}+185\mathbf{t}^{12}+74\mathbf{t}^{14}$ & $21\mathbf{t}^{5}+60\mathbf{t}^{7}+116\mathbf{t}^{9}+153\mathbf{t}^{11}+137\mathbf{t}^{13}+65\mathbf{t}^{15}$ & $10\mathbf{t}^{10}+24\mathbf{t}^{12}+31\mathbf{t}^{14}+20\mathbf{t}^{16}$ & $\mathbf{t}^{17}$ \\
3 & $56\mathbf{t}^{2}+140\mathbf{t}^{4}+210\mathbf{t}^{6}+220\mathbf{t}^{8}+154\mathbf{t}^{10}+56\mathbf{t}^{12}$ & $56\mathbf{t}^{3}+175\mathbf{t}^{5}+310\mathbf{t}^{7}+374\mathbf{t}^{9}+296\mathbf{t}^{11}+121\mathbf{t}^{13}$ & $35\mathbf{t}^{6}+100\mathbf{t}^{8}+164\mathbf{t}^{10}+162\mathbf{t}^{12}+80\mathbf{t}^{14}$ & $10\mathbf{t}^{11}+20\mathbf{t}^{13}+15\mathbf{t}^{15}$ & $0$ \\
4 & $70\mathbf{t}^{3}+175\mathbf{t}^{5}+230\mathbf{t}^{7}+181\mathbf{t}^{9}+70\mathbf{t}^{11}$ & $70\mathbf{t}^{4}+210\mathbf{t}^{6}+320\mathbf{t}^{8}+287\mathbf{t}^{10}+125\mathbf{t}^{12}$ & $35\mathbf{t}^{7}+90\mathbf{t}^{9}+111\mathbf{t}^{11}+61\mathbf{t}^{13}$ & $5\mathbf{t}^{12}+6\mathbf{t}^{14}$ & $0$ \\
5 & $56\mathbf{t}^{4}+126\mathbf{t}^{6}+127\mathbf{t}^{8}+56\mathbf{t}^{10}$ & $56\mathbf{t}^{5}+147\mathbf{t}^{7}+169\mathbf{t}^{9}+84\mathbf{t}^{11}$ & $21\mathbf{t}^{8}+42\mathbf{t}^{10}+29\mathbf{t}^{12}$ & $\mathbf{t}^{13}$ & $0$ \\
6 & $28\mathbf{t}^{5}+49\mathbf{t}^{7}+28\mathbf{t}^{9}$ & $28\mathbf{t}^{6}+56\mathbf{t}^{8}+36\mathbf{t}^{10}$ & $7\mathbf{t}^{9}+8\mathbf{t}^{11}$ & $0$ & $0$ \\
7 & $8\mathbf{t}^{6}+8\mathbf{t}^{8}$ & $8\mathbf{t}^{7}+9\mathbf{t}^{9}$ & $\mathbf{t}^{10}$ & $0$ & $0$ \\
8 & $\mathbf{t}^{7}$ & $\mathbf{t}^{8}$ & $0$ & $0$ & $0$ \\
\hline
\end{tabular}

\end{center}
各 $(g,\beta)$ について $\sum_{2J_r}\hN_{[1],g,\beta,J_r}(\mathrm{K1})\,\bt^{2J_r}$．$g$ の最大値 8 は K1 の種数と一致する．
\end{landscape}

\section{データと再現性}

計算に用いたコード，結果（係数辞書と表），自動テストは，研究用リポジトリ（ローカル）にある．
\begin{itemize}
\item コード：\texttt{src/rcsknots/}
\item refined CS 不変量：\texttt{data/cables/}
\item BPS 不変量：\texttt{data/bps/}
\item 規約の詳細：\texttt{docs/conventions.md}
\item 負の係数の原因の切り分け（\S\ref{sec:cause}）：\\ \texttt{scripts/negativity\_cause.py}，\texttt{scripts/framing\_scan.py}，\\ 結果は \texttt{data/negativity\_cause.md}
\item 圏化との比較（\S\ref{sec:categorification}）：\\ \texttt{scripts/homology\_invariants.py}，\texttt{scripts/compare\_khovanov.py}，\\ \texttt{scripts/construct\_superpolynomial.py}，\texttt{scripts/superpoly\_lmov.py}，\\ 結果は \texttt{data/homology/}，\texttt{data/superpoly/}
\end{itemize}

%======================================================================
\begin{thebibliography}{99}
\bibitem{KN} S.~Kameyama and S.~Nawata, \emph{Refined large N duality for knots}, J.~Knot Theory Ramifications \textbf{32} (2023) 2041001, arXiv:1703.05408.
\bibitem{AS} M.~Aganagic and S.~Shakirov, \emph{Knot homology from refined Chern--Simons theory}, Commun.~Math.~Phys.\ \textbf{333} (2015) 187, arXiv:1105.5117.
\bibitem{AS2} M.~Aganagic and S.~Shakirov, \emph{Refined Chern--Simons theory and knot homology}, Proc.~Symp.~Pure Math.\ \textbf{85} (2012) 3, arXiv:1202.2489.
\bibitem{Ch} I.~Cherednik, \emph{Jones polynomials of torus knots via DAHA}, arXiv:1111.6195.
\bibitem{GN} E.~Gorsky and A.~Neguţ, \emph{Refined knot invariants and Hilbert schemes}, J.~Math.~Pures Appl.\ \textbf{104} (2015) 403, arXiv:1304.3328.
\bibitem{CD14} I.~Cherednik and I.~Danilenko, \emph{DAHA and iterated torus knots}, Algebr.~Geom.~Topol.\ \textbf{16} (2016) 843, arXiv:1408.4348.
\bibitem{CD15} I.~Cherednik and I.~Danilenko, \emph{DAHA approach to iterated torus links}, arXiv:1509.08351.
\bibitem{CP} I.~Cherednik and I.~Philipp, \emph{DAHA and plane curve singularities}, Algebr.~Geom.~Topol.\ \textbf{18} (2018) 333, arXiv:1605.00978.
\bibitem{MS} H.~Morton and P.~Samuelson, \emph{The HOMFLYPT skein algebra of the torus and the elliptic Hall algebra}, Duke Math.~J.\ \textbf{166} (2017) 801, arXiv:1410.0859.
\bibitem{GS} S.~Gukov and M.~Stošić, arXiv:1112.0030（題名・掲載誌未確認）.
\bibitem{DGR} N.~Dunfield, S.~Gukov and J.~Rasmussen, \emph{The superpolynomial for knot homologies}, Experiment.~Math.\ \textbf{15} (2006) 129, arXiv:math/0505662.
\bibitem{DMMSS} P.~Dunin-Barkowski, A.~Mironov, A.~Morozov, A.~Sleptsov and A.~Smirnov, \emph{Superpolynomials for toric knots from evolution induced by cut-and-join operators}, JHEP \textbf{03} (2013) 021, arXiv:1106.4305.
\bibitem{MMSS} A.~Mironov, A.~Morozov, Sh.~Shakirov and A.~Sleptsov, \emph{Interplay between MacDonald and Hall--Littlewood expansions of extended torus superpolynomials}, JHEP \textbf{05} (2012) 070, arXiv:1201.3339.
\bibitem{Dan} I.~Danilenko, arXiv:1405.0884（題名未確認）.
\bibitem{Mo18} A.~Morozov, Phys.~Lett.~B \textbf{782} (2018), arXiv:1801.02407（題名未確認）.
\bibitem{AnMo} An.~Morozov, \emph{The first-order deviation of superpolynomial in an arbitrary representation from the special polynomial}, JETP Lett.\ \textbf{97} (2013) 171, arXiv:1211.4596.
\bibitem{KR2} M.~Khovanov and L.~Rozansky, \emph{Matrix factorizations and link homology II}, arXiv:math/0505056.
\bibitem{Kh07} M.~Khovanov, \emph{Triply-graded link homology and Hochschild homology of Soergel bimodules}, arXiv:math/0510265.
\bibitem{Ras} J.~Rasmussen, \emph{Some differentials on Khovanov--Rozansky homology}, arXiv:math/0607544.
\bibitem{ORS} A.~Oblomkov, J.~Rasmussen and V.~Shende, \emph{The Hilbert scheme of a plane curve singularity and the HOMFLY homology of its link}, arXiv:1201.2115.
\bibitem{HM} M.~Hogancamp and A.~Mellit, \emph{Torus link homology}, arXiv:1909.00418.
\bibitem{NS} K.~Nakagane and T.~Sano, \emph{Computations of HOMFLY homology}, arXiv:2111.00388.
\bibitem{GHM} E.~Gorsky, M.~Hogancamp and A.~Mellit, \emph{Tautological classes and symmetry in Khovanov--Rozansky homology}, Duke Math.~J.\ \textbf{173} (2024) 2481, arXiv:2103.01212.
\bibitem{GGS} E.~Gorsky, S.~Gukov and M.~Stošić, \emph{Quadruply-graded colored homology of knots}, arXiv:1304.3481.
\bibitem{khoca} L.~Lewark, \texttt{khoca}（Python パッケージ，版 1.5 を使用）．
\bibitem{snappy} SnapPy（Python パッケージ，版 3.3.2 を使用．knot Floer ホモロジーは \texttt{knot\_floer\_homology} 1.2.2）．
\bibitem{GL} P.~Galashin and T.~Lam, \emph{Monotone links in DAHA and EHA}, arXiv:2307.16794.
\end{thebibliography}

\end{document}
